
\documentclass[a4paper,12pt]{report} 

\usepackage[top = 3.4cm, left = 2.4cm, right = 2.4cm, bottom = 2.6cm]{geometry}
\usepackage{multicol}

\usepackage{tikz}
\usepackage{draftwatermark}
\SetWatermarkText{
\begin{tikzpicture}
\draw[fill] (-1.4142, 0) -- (0, -1) -- (1.4142, 0) -- (0, 1) -- cycle;
\end{tikzpicture}
}
\SetWatermarkHorCenter{2cm}
\SetWatermarkVerCenter{2cm}
\SetWatermarkColor[rgb]{0.35,0.65,0.9}

\usepackage{hyperref}
\hypersetup{
	colorlinks=true,
	urlcolor=black
}

\usepackage{natbib}
%\usepackage{apacite}
%\bibliographystyle{apacite}
\usepackage{bibentry}
\nobibliography*

\begin{document} 

\begin{center}
{\Large \textbf{Lyron Winderbaum}} \\ \vspace{4mm}
\end{center}

\section*{Professional Knowledge}

Central to my teaching ideology is that each student has individual learning needs, and that the best most practical way to address these is use student voice and student driven content to inform my teaching. I build relationships and rapport with students through one-on-one conversations to create an environment in which they feel safe and comfortable communicating with me honestly. This allows for ongoing informal formative assessment of their learning, which I use to differentiate and adjust my teaching to their needs. I offer choices to my students in my classrooms as much as possible, as I beleive this provides opportunities for students to differentiate their own learning by choosing an option more beneficial to them as individuals, while also supporting and encouraging them to be self-reflective about how they learn.

I beleive that when students make choices about their own learning, they can become more invested in the activity, be more engaged, and get more out of it. To this end I encourage students to take an active role in shaping their learning in our shared classrooms. In practice this takes many different forms: students asking for me to speed up or slow down as we go through subject matter, students engaging is self-directed investigations, to customisation of the physical and social environment through negotiation and guided self-reflection. The depth of my knowledge in mathematics and the sciences allows me to easily extend students well beyond their year level (or even beyond school level curriculum as appropriate), give context to the learning by placing it in an authentic framework. My experience working at the maths learning centre has exposed me to many differnt ways to understand the same mathematical concepts, which can be a very valuble when students are struggling to understand something. Beyond simple extension and alterative explanations of concepts, my experience in mathematics and science also allow me to guide students on self-directed investigations of great depth, which I thoroughly enjoy. For example, if doing a practical in science and students suggest changing the experimental design in some way, I will guide them to make adjustments that make sense, and do everything in my power to support this happening.

\section*{Professional Practice}

For students to learn and excel they need to feel comfortable and safe. As such, an affective approach to student wellbeing is crucial. By producing an environment in which students can feel safe, I can help them to discover their strengths and excel. I do this primarily by building relationships with my students. This is a subtle process that takes some time, but gradually trust is developed and through that I can gain some insight into how my students are feeling and make adjustments to the classroom environment in order to help them feel more comfortable. For example, in one of my classes I ran an activity in which I taught the class how to fold A4 paper into silver rhombic dodecahedra and had them investigate the properties of this shape. The process of folding the paper is quite therapeutic, and as I ran this activity after already having had the class for a few weeks and having established relationships with them enough that they felt comfortable coming up to me and telling about their feelings, one student came up to me after class and admitted that she had been suffering from severe anxiety and that had been the reason that she had missed alot of classes but that folding the paper rhombic units had actually helped her feel a lot more comfortable in class. As a result I modified the classroom environment the next week by putting a stack of paper in the corner of the class and informing the class that if anyone was feeling overwhelmed or anxious during class that they could move over into the corner and fold some rhombic units for dodecahedra as a calming exercise. Several students beyond the original one who had come to talk to me took advantage of this corner in the weeks to follow, and several of them told me about how it had helped them calm down after having had some pretty severely traumatic events outside of the class. Some of these students attendance rates also improved after this. Its in ways like this that I aim to create a safe environment in which my students can feel comfortable and productive. If students are in an uncomfortable environment or feel unsafe it is completely unreasonable to expect them to be able to learn, so as teachers providing a safe environment for our students has to be one of our highest priorities, not just because without that it would be impossible for our students to learn, but because everybody deserves to feel safe.

\section*{Professional Engagement}

Learning is a passion of mine, and I have never stopped learning at any point during my life. Through discussion with colleagues and self-reflection I identify my weaknesses and areas of ignorance, and constantly work to better myself and improve. I particularly benefit from open discussions with colleagues who hold very different views/ beliefs and use very different teaching methods to me, as this is when I learn the most. For example on my placement to BHS, one of my mentor teachers had a very different teaching style to mine. Through discussions with her and with the students I gained alot of insight into how there truely is no single best teaching style, there is simply teaching styles that better suit certain students, and teaching styles that better suit different teachers. I learnt that there is alot of value in learning from teachers with different teaching styles than your own, because pushing yourself outside of your comfort zone and learning these different skills can be incredibly valuable with faced with students who do not respond to your preferred style. Professional development is always a pleasure for me, and while I was at BHS and they had a numeracy focused professional development day I offered to take a workshop on probabiltiy and share my expertise with the rest of the staff. I even had one staff member approach me several days later and thank me, because she had apparently suffered from severe maths anxiety, having a minor panic attack on the morning of the PD day because she realised it would have a numeracy focus, but she said that the workshop I ran had opened her eyes to the possibility that maths might not actually be all that scary and more connected to concepts with which she felt comfortable than she had previously realised. I was very proud of her, and I only wish I could bring experiences like that to more people.

\end{document}
